% Source: https://artofproblemsolving.com/community/c3h1443927_x_and_y
% Stable identity: aops-topic-1443927-q01
Clearly $x>0$ and $y>0$ (if $x=0$ the first equation fails; if $y=0$ the second fails; and both square roots need non-negative arguments). Divide:
\[
1+\dfrac{1}{x+y}=\dfrac{2}{\sqrt{3x}},\qquad 1-\dfrac{1}{x+y}=\dfrac{4\sqrt2}{\sqrt{7y}} .
\]
Adding and subtracting,
\[
1=\dfrac{1}{\sqrt{3x}}+\dfrac{2\sqrt2}{\sqrt{7y}},\qquad \dfrac{1}{x+y}=\dfrac{1}{\sqrt{3x}}-\dfrac{2\sqrt2}{\sqrt{7y}} .
\]
Multiplying these two equations,
\[
\dfrac{1}{x+y}=\dfrac{1}{3x}-\dfrac{8}{7y}\quad\Longrightarrow\quad 21xy=(7y-24x)(x+y)\quad\Longrightarrow\quad 7y^2-38xy-24x^2=0 .
\]
This factors as $(7y+4x)(y-6x)=0$, and since $x,y>0$ we get $y=6x$. Then $\dfrac{2\sqrt2}{\sqrt{7y}}=\dfrac{2\sqrt2}{\sqrt{42x}}=\dfrac{2}{\sqrt{21x}}$, and the first equation becomes
\[
\dfrac{1}{\sqrt x}\left(\dfrac{1}{\sqrt3}+\dfrac{2}{\sqrt{21}}\right)=1\quad\Longrightarrow\quad x=\left(\dfrac{1}{\sqrt3}+\dfrac{2}{\sqrt{21}}\right)^2=\dfrac{1}{3}+\dfrac{4}{21}+\dfrac{4}{\sqrt{63}}=\dfrac{11+4\sqrt7}{21}.
\]
So the unique solution is
\[
\boxed{x=\dfrac{11+4\sqrt7}{21},\qquad y=6x=\dfrac{22+8\sqrt7}{7}} .
\]
Conversely these values satisfy both derived equations (sum and difference), hence the original ones; this is confirmed numerically.
